\documentclass[11pt]{article}
\usepackage[margin=1in]{geometry}
\usepackage{amsmath,amssymb,amsthm}
\usepackage{hyperref}
\hypersetup{hidelinks}
\usepackage{booktabs}

\newtheorem{theorem}{Theorem}
\newtheorem{lemma}[theorem]{Lemma}
\newtheorem{corollary}[theorem]{Corollary}
\theoremstyle{definition}
\newtheorem{definition}{Definition}

\setlength{\parskip}{4pt plus 1pt minus 1pt}
\setlength{\parindent}{0pt}

\title{Escalation Dominates Top-Down Allocation\\Under Bounded Context}
\author{Patrick D.\ McCarthy}
\date{}

\begin{document}
\maketitle

\begin{abstract}
Top-down allocation --- where a supervisory process directs
subproblems to lower levels --- is viable when the action repertoire
$|A|$ fits within bounded active context $C_n$. We prove that once
$|A|$ exceeds a crossover point $|A|^* = C_n / 2\alpha$, top-down
allocation degrades inference quality monotonically, whether the
architecture is flat ($O(|A|)$ overhead) or hierarchical ($O(b \cdot
\log_b |A|)$ overhead for branching factor $b$). The crossover is
permanent: recovery requires growing $C_n$ with $|A|$, which no
bounded system can satisfy. Escalation --- bottom-up failure signaling
where higher levels engage only when lower levels fail --- achieves
$O(1)$ steady-state context cost and \emph{improves} with capability
growth, as successfully encoded actions relieve rather than burden
the highest level. Since intelligence requires $|A| \gg C_n$, every
intelligent system operates beyond the crossover. Escalation dominance
is not a universal claim but a scaling result: top-down wins early,
escalation wins at every scale that matters. This is a formal critique
of precision allocation models including the Free Energy Principle
\cite{Friston2010,Friston2017} and the System~M routing architecture
of Dupoux, LeCun, and Malik \cite{DupouxLeCunMalik2026}.
\end{abstract}

% ------------------------------------------------------------------
\section{Introduction}
\label{sec:intro}
% ------------------------------------------------------------------

Every system that acts under uncertainty must solve a control problem:
which problems receive active reasoning, and which are handled by
automatic, pre-encoded responses? This question is distinct from the
content of reasoning; it concerns the \emph{architecture of
attention}. Under bounded active context $C_n$ --- a hard limit on the
number of propositions simultaneously available for inference --- the
control architecture is not a free parameter but a survival-critical
design choice.

Two competing architectures present themselves:

\begin{enumerate}
\item \textbf{Top-down allocation.} A supervisory process at the
   highest encoding level $L_n$ monitors all subsystems, evaluates
   relevance, and routes problems to appropriate levels. This
   architecture is implicit in precision weighting under the Free
   Energy Principle \cite{Friston2010,Friston2017}, where higher
   cortical levels set the gain on prediction errors at lower levels.
   It is explicit in the System~M proposal of Dupoux, LeCun, and Malik
   \cite{DupouxLeCunMalik2026}, where a meta-controller directs when
   Systems~A (observation) and B (action) engage. Kahneman's popular
   account of System~1 and System~2 \cite{Kahneman2011} implicitly
   assumes top-down control: System~2 ``decides'' when to override
   System~1.

\item \textbf{Escalation.} Lower levels handle problems autonomously.
   Higher levels engage only when lower levels fail --- that is, when
   uncertainty cannot be bounded at the current level. This
   architecture is implicit in Global Workspace Theory
   \cite{Baars1988,Dehaene2011}, where local processors broadcast to
   the global workspace only when local resolution fails.
\end{enumerate}

At small scale --- when $|A| \leq C_n$, so the entire action repertoire
fits in active context --- top-down allocation is viable and may even be
preferable: the supervisor can hold complete state with room to spare,
and the benefit of optimal routing outweighs the modest overhead.
Intelligence, however, requires $|A| \gg C_n$: an entity must know far
more than it can hold in mind at once. This paper proves that once
$|A|$ exceeds $C_n$, top-down allocation degrades and escalation
dominates --- and that the crossover is permanent. The result is not
that top-down is wrong but that it cannot scale, and scaling is what
intelligence demands.

The proof proceeds as follows. Section~2 introduces the encoding
hierarchy $E(c,r)$ and the formal definitions required for the result.
Section~3 establishes the escalation principle (Theorem~1): informed
actions are handled at the lowest sufficient level. Section~4 proves
the central result (Theorem~2): top-down allocation degrades inference
quality as capability grows, with recovery requiring $C_n$ to grow
linearly with capability --- a condition no bounded system can satisfy.
Section~5 compares the two architectures directly. Section~6 engages
with prior work.

This paper is the second in a series. Paper~1 \cite{McCarthy2026a}
established that any system accumulating knowledge under bounded
context necessarily maintains graph structure. The present paper
addresses how that structure is \emph{controlled} --- how problems are
routed through the encoding hierarchy.

% ------------------------------------------------------------------
\section{Formal Framework}
\label{sec:framework}
% ------------------------------------------------------------------

We introduce the minimal definitions required for the main results.
We inherit the world state space $W$, proposition, knowledge base $K$,
uncertainty $U(w,K) = -\log P(w \mid K)$ \cite{Cover2006}, and bounded
active context $C_n$ from \cite{McCarthy2026a}. An \emph{informed
action} $a \in A$ is an action grounded in propositions from $K$; we
write $K^{[a]} \subseteq K$ for the set of propositions that ground $a$.

\begin{definition}[Encoding Hierarchy]
\label{def:hierarchy}
An \emph{encoding hierarchy} is a continuous two-dimensional space
$E(c,r)$ parameterized by:
\begin{itemize}
\item \textbf{Runtime cost} $c \geq 0$: the computational cost of
   accessing and modifying the encoding.
\item \textbf{Reversibility} $r \in [0,1]$: the ease with which the
   encoding can be revised. $r = 1$ is fully reversible; $r = 0$ is
   permanent.
\end{itemize}

Boundaries within this continuous space are \emph{interfaces}
$L_0, L_1, \ldots, L_n$. The number of interfaces and their positions
are substrate-determined: biological systems, artificial systems, and
institutional systems partition $E(c,r)$ differently.
\end{definition}

Table~\ref{tab:hierarchy} illustrates the hierarchy with representative
interfaces. The precise number and character of levels varies by
substrate; what is invariant is the monotonic relationship between
runtime cost and reversibility.

\begin{table}[h]
\centering
\small
\begin{tabular}{@{}llll@{}}
\toprule
\textbf{Interface} & \textbf{Description} & \textbf{Runtime Cost} &
\textbf{Reversibility} \\
\midrule
$L_n$ & Active reasoning / working memory & $O(\text{compute})$ &
Fully reversible ($r \approx 1$) \\
$L_{n-1}$ & Learned heuristics / cached policies & $O(1)$ retrieval &
Revisable with effort \\
$\vdots$ & \emph{substrate-dependent intermediate levels} & & \\
$L_1$ & Developmental encoding / early learning & Near-zero &
Difficult to revise \\
$L_0$ & Heritable / constitutional encoding & $0$ &
Permanent ($r \approx 0$) \\
\bottomrule
\end{tabular}
\caption{Representative encoding hierarchy. Cost increases and
reversibility increases from $L_0$ to $L_n$. The number of interfaces
is substrate-determined.}
\label{tab:hierarchy}
\end{table}

\begin{definition}[Active Context Capacity and Relevance Density]
\label{def:context}
The \emph{active context capacity} $C_n$ is bounded: $|\mathrm{ctx}(L_n,
t)| \leq C_n$. For an informed action $a$, the \emph{relevance
density} at $L_n$ is:
\[
\rho(L_n, f) = \frac{|K^{[a]}|}{|\mathrm{ctx}(L_n, t)|}
\]
where $K^{[a]} \subseteq K$ is the knowledge relevant to $a$.
$\rho = 1$ is maximum signal: active context holds exactly what $a$
requires. $\rho \to 0$ is noise-dominated: context is filled with
content irrelevant to $a$.
\end{definition}

\begin{definition}[Delegation Overhead]
\label{def:delegation}
Under \emph{top-down allocation}, $L_n$ actively manages lower-level
processes. For each active delegation, $L_n$ must hold in context:
\begin{enumerate}
\item[(i)] The \emph{subproblem specification}: what was delegated.
\item[(ii)] The \emph{expected return signal}: what a successful
   result looks like.
\item[(iii)] The \emph{integration context}: how the result will be
   incorporated into the current reasoning state.
\end{enumerate}
Let $\alpha > 0$ be the \emph{minimum context cost per delegation}:
the irreducible number of context slots consumed by maintaining one
active delegation. $\alpha > 0$ because each delegation requires at
least one proposition specifying what was delegated.
\end{definition}

% ------------------------------------------------------------------
\section{Escalation}
\label{sec:escalation}
% ------------------------------------------------------------------

\begin{theorem}[Escalation]
\label{thm:escalation}
Under the encoding hierarchy (Definition~\ref{def:hierarchy}), the
cost-minimizing control policy handles each action $a$ at the lowest
level $L_i$ such that the knowledge encoded at $L_i$ is sufficient to
bound uncertainty for $a$. Escalation to $L_{i+1}$ occurs if and only
if $L_i$ fails to bound uncertainty.
\end{theorem}

\begin{proof}
By Definition~\ref{def:hierarchy}, higher levels have strictly greater
runtime cost: $c(L_{i+1}) > c(L_i)$. This monotonicity is not
incidental but reflects the architecture of complexity: Simon
\cite{Simon1962} shows that stable complex systems are
\emph{nearly decomposable} --- they process information locally, with
infrequent communication between subsystems. The encoding hierarchy
formalizes this structure: lower levels are cheaper precisely because
they are more specialized and less revisable. We show that handling $a$
at any level higher than the lowest sufficient level is strictly
dominated.

Let $L_i$ be the lowest level such that $U(w, K^{[a]}_{L_i}) \leq
\epsilon$ for the action-relevant threshold $\epsilon$. Consider the
alternative of handling $a$ at $L_j$ ($j > i$):

\begin{enumerate}
\item \textbf{Cost.} $c(L_j) > c(L_i)$ by the monotonicity of cost in
   the hierarchy. The uncertainty reduction is the same (both levels
   achieve $U \leq \epsilon$), so the excess cost $c(L_j) - c(L_i) > 0$
   yields zero marginal benefit. The higher-level policy is strictly
   dominated on cost.

\item \textbf{Context.} If $j = n$ (the active reasoning level), then
   handling $a$ at $L_n$ consumes active context capacity $C_n$. Since
   $C_n$ is bounded and $L_n$ is the only level that handles genuinely
   novel problems (problems no lower level can resolve), each unnecessary
   use of $L_n$ reduces $\rho$ for the next novel problem. The
   higher-level policy is strictly dominated on context conservation.
\end{enumerate}

Under any cost-minimizing criterion --- whether natural selection,
engineering optimization, or resource-rational design --- the dominated
policy is eliminated. The unique cost-minimizing policy is: handle $a$
at the lowest sufficient level; escalate if and only if that level
fails.

\textbf{Remark.} This result does not claim that all systems
\emph{implement} escalation. Path dependence, architectural
constraints, and design choices can produce systems that use higher
levels than necessary. The theorem establishes that such systems pay a
strictly positive cost premium: they are not cost-optimal under bounded
context. The claim is normative (escalation is optimal) rather than
descriptive (all systems escalate).
\end{proof}

\begin{corollary}[Context Window as Frontier]
\label{cor:frontier}
$L_n$ is the frontier of unbounded uncertainty, not the seat of
intelligence. A problem reaching $L_n$ is evidence of failure at every
lower level. An entity that routes most problems through $L_n$ is not
``thinking more'' --- it is failing to encode.
\end{corollary}

\begin{corollary}[Simultaneous Multi-Level Activity]
\label{cor:simultaneous}
All levels of the hierarchy operate concurrently. $L_0$ handles
heritable responses, $L_1$ handles developmental encodings, and
intermediate levels handle their respective domains ---
simultaneously. $L_n$ handles only the residual: problems that no
lower level can resolve. The escalation architecture does not
serialize processing through a hierarchy; it parallelizes across levels,
with upward propagation occurring only on failure.
\end{corollary}

% ------------------------------------------------------------------
\section{Top-Down Allocation Degrades Inference Quality}
\label{sec:topdown}
% ------------------------------------------------------------------

\begin{theorem}[Top-Down Allocation Degrades Inference Quality]
\label{thm:topdown}
Under top-down allocation, the supervisory level $L_n$ must maintain
delegation state for managed subprocesses
(Definition~\ref{def:delegation}). Let $k(t)$ be the number of
concurrent delegations at time $t$. The relevance density for a novel
problem $f_{\mathrm{novel}}$ is:
\[
\rho(L_n, f_{\mathrm{novel}}) =
\frac{|K^{[f_{\mathrm{novel}}]}|}{|K^{[f_{\mathrm{novel}}]}| + k(t)
\cdot \alpha}
\]
If $k(t)$ grows with $|A|$, then $\rho \to 0$ as $|A| \to \infty$.
Recovery requires growing $C_n$ with $k(t)$ --- a condition no bounded
system can satisfy.
\end{theorem}

\begin{proof}
Under top-down allocation, $L_n$ serves two simultaneous functions:
(a) reasoning about the novel problem $f_{\mathrm{novel}}$, requiring
$|K^{[f_{\mathrm{novel}}]}|$ context slots, and (b) maintaining
delegation state for $k(t)$ concurrently managed subprocesses, each
consuming $\alpha$ context slots (Definition~\ref{def:delegation}).

The total context consumed is:
\[
|\mathrm{ctx}(L_n, t)| = |K^{[f_{\mathrm{novel}}]}| + k(t) \cdot \alpha
\]

Since $\alpha > 0$, any growth in $k(t)$ increases the denominator
while the numerator $|K^{[f_{\mathrm{novel}}]}|$ is bounded per
problem. Therefore $\rho$ decreases monotonically with $k(t)$.

The critical question is whether $k(t)$ grows with $|A|$. We show that
it must under any architecture that merits the name ``top-down'':

\textbf{Why $k$ grows with $|A|$.} Top-down allocation is defined by
the supervisory level \emph{initiating} and \emph{tracking}
delegations (Definition~\ref{def:delegation}). A supervisor that
delegates but does not track (fire-and-forget) holds no return-integration context and
cannot evaluate, redirect, or integrate results --- it is not
supervising. A supervisor that tracks only a fixed subset of $|A|$
while ignoring the rest has abandoned top-down control of the
ignored actions, which are then self-managing (i.e., operating under
escalation). In either case, the system is no longer top-down for the
untracked actions.

For the portion of $A$ that remains under genuine top-down
supervision, $k(t)$ grows with the number of supervised actions. Under
flat top-down, $k = |A_{\text{supervised}}|$. Under hierarchical
top-down with branching factor $b$, each supervisory level tracks $b$
subordinates across $\lceil \log_b |A_{\text{supervised}}| \rceil$
levels, giving $k = b \cdot \log_b |A_{\text{supervised}}|$. In both
cases $k$ grows with the supervised repertoire.

The critical threshold occurs when $k(t) \cdot \alpha \geq C_n -
|K^{[f_{\mathrm{novel}}]}|$: delegation overhead has consumed all
context capacity not occupied by the current problem's knowledge.

To maintain $\rho$ at a fixed level as $|A|$ grows, $C_n$ must grow
at least linearly with $k(t)$. But $C_n$
is bounded by substrate (Definition~\ref{def:context}). Therefore
top-down allocation necessarily degrades inference quality as the
supervised repertoire grows under bounded context.
\end{proof}

\subsection{The crossover: when top-down stops working}
\label{sec:crossover}

Top-down allocation is viable when delegation overhead fits comfortably
within $C_n$. Define the \emph{crossover point} as the supervised
repertoire size at which delegation overhead consumes half of active
context:
\[
k(|A|^*) \cdot \alpha = \frac{C_n}{2}
\]
For flat top-down ($k = |A|$), this gives $|A|^* = C_n / 2\alpha$. For
hierarchical top-down ($k = b \cdot \log_b |A|$), the crossover is
$b \cdot \log_b |A|^* = C_n / 2\alpha$, which gives a larger $|A|^*$
--- hierarchy delays the crossover but does not eliminate it, because
$\log_b |A|$ still grows without bound.

Below $|A|^*$, the supervisor retains at least $C_n/2$ slots for
reasoning --- adequate for most problems. Above $|A|^*$, delegation
overhead dominates: every new supervised capability \emph{reduces} the
system's ability to reason about novel problems.

\textbf{Estimating the crossover.} The minimum delegation cost $\alpha$
is at least 1: specifying ``what was delegated'' requires at least one
proposition in context. A tighter estimate: each delegation requires a
specification, an expected return, and an integration slot
(Definition~\ref{def:delegation}), giving $\alpha \geq 3$. Under
Cowan's \cite{Cowan2001} revised estimate of working memory capacity
($C_n \approx 4$, not Miller's $7 \pm 2$), the flat crossover is
$|A|^* = C_n / 2\alpha \approx 4/6 < 1$. Even with Miller's more
generous bound ($C_n \approx 7$), $|A|^* \approx 7/6 \approx 1$.
The crossover for biological working memory occurs at \emph{one
delegation} --- a supervisor that manages even a single subprocess has
already consumed most of its reasoning capacity. This is why human
``multitasking'' degrades performance: the delegation overhead of
tracking two tasks consumes the context needed to reason about
either.

This crossover is permanent under any top-down architecture. Once
$k(|A|) \cdot \alpha > C_n / 2$, the only recovery is to grow $C_n$
--- which is bounded by substrate --- or to abandon top-down
supervision for part of $A$, which is to adopt escalation for those
actions.

The crossover clarifies the scope of the result. Top-down allocation
is not wrong --- it is a \emph{small-scale} architecture. For an entity
whose entire repertoire fits in active context, top-down supervision is
efficient and may be optimal. But intelligence requires $|A| \gg C_n$:
an entity that can only act on what it can simultaneously hold in mind
is not intelligent but reactive. Every physical system that we
recognize as intelligent --- from insects routing foraging with stigmergic
signals to humans operating with $\sim$7 active context slots across
thousands of learned skills --- operates well beyond the crossover.

\textbf{Cell division as illustration.} A single-celled organism can
``manage'' its repertoire top-down: the entire state fits in one
context. A multicellular organism with $10^{13}$ cells cannot manage
each cell individually through a central controller.

Biology's actual solution is a \emph{hybrid}: local autonomy (cells
monitor nutrient gradients, contact inhibition, and growth factor
concentrations) combined with systemic signaling (hormonal regulation,
tumor suppressors, immune surveillance). The systemic signals are
genuine top-down control --- but they are \emph{broadcast} signals
(growth hormone affects all cells of the right type), not individual
delegations (no central process tracks each cell's state). Broadcasting
is $O(1)$ in delegation context: the signal is sent once regardless of
the number of recipients. Individual delegation would require tracking
$10^{13}$ cells, which is well beyond any bounded context.

The hybrid architecture is precisely the crossover result: top-down
broadcast control for system-wide coordination (below the crossover for
broadcast), combined with local escalation for individual cell
decisions (above the crossover for per-cell management). When the
systemic signaling fails --- when a cell ignores growth suppression
signals --- the immune system engages: escalation on failure.

\subsection{The hierarchical top-down objection}

The strongest objection to Theorem~\ref{thm:topdown} is that it attacks
flat top-down allocation --- a strawman. Real supervisory systems are
\emph{hierarchical}: a top level manages $b$ intermediaries, each
managing $b$ sub-intermediaries, so that overhead at any single level is
$O(b)$, not $O(|A|)$. Does hierarchical decomposition rescue top-down
allocation?

No. The hierarchical decomposition is more subtle than the flat case,
but it does not rescue top-down allocation. We address the strongest
version of the objection directly.

\textbf{(1) $L_n$ sees only $O(b)$ overhead --- but the hierarchy has
a different cost.}
The strongest version of the hierarchical objection notes that $L_n$
manages only $b$ direct subordinates, so overhead at $L_n$ is
$b \cdot \alpha$ --- constant in $|A|$. The $\log_b |A|$ levels of
overhead are distributed across intermediate levels, each with its own
bounded context. No single level overflows. This is correct and is a
genuine advantage of hierarchical decomposition over flat top-down.

However, the cost is not eliminated --- it is \emph{redistributed}.
The hierarchy requires $\lceil \log_b |A| \rceil$ supervisory levels,
each of which devotes $b \cdot \alpha$ of its context to delegation
state. The number of management-dedicated levels grows with $|A|$.

\textbf{(2) Management levels displace knowledge levels.}
This is the deeper cost, and it has a precise analogue in organizational
economics. Williamson \cite{Williamson1967} proves that each management
layer in a hierarchical organization introduces ``control loss'': a
fraction $\alpha_W < 1$ of the intended directive is lost at each level,
so that $n$ management levels transmit only $\alpha_W^n$ of the
top-level intent. This limits the viable depth of hierarchical control
and, consequently, organization size. The same structure applies here:
each intermediary in a hierarchical top-down system is itself a
supervisory process that must hold delegation context, and the
cumulative cost of these management layers limits the encoding
capacity available for knowledge.

In an escalation architecture, every encoding level is a
\emph{knowledge level}: it holds proven actions and propositions, not
routing state. In hierarchical top-down, $\lceil \log_b |A| \rceil$
levels are \emph{management levels}: they hold delegation state. As
$|A|$ grows, the number of management levels grows (logarithmically),
consuming encoding capacity that escalation would use for knowledge.
Williamson's result shows this is not merely wasteful but structurally
limiting: the management overhead places an upper bound on the
effective depth of hierarchical control, just as it places an upper
bound on firm size.

The cost is not context overflow at any single level but a growing
fraction of the entity's encoding hierarchy devoted to management rather
than knowledge. An entity with 10 encoding levels and $|A|$ requiring
4 management levels has only 6 levels for knowledge; escalation would
use all 10. The management overhead is in \emph{levels consumed}, not
\emph{context slots consumed at $L_n$}.

\textbf{(3) The causal direction is the fundamental difference.}
The deepest distinction between hierarchical top-down and escalation is
not the cost asymptotic but the \emph{causal direction of engagement}:

\begin{itemize}
\item In hierarchical top-down, the upper level \emph{initiated} the
   delegation. It is waiting for a result it requested. It must hold
   return-integration context: what was delegated, what a successful
   result looks like, and how to incorporate it. This context is held
   \emph{whether or not the subprocess encounters difficulty}. A
   supervisor waiting for ten reports holds ten slots regardless of
   whether any report is late.

\item In escalation, the upper level \emph{did not initiate} anything.
   It holds zero context for problems handled below. It loads context
   only when a failure signal arrives from below --- and in steady
   state, failures are rare because lower levels handle most problems.
   There is nothing to ``wait for'' because nothing was requested.
\end{itemize}

This is not a quantitative difference in overhead --- it is a structural
difference in what context is held and why. Hierarchical top-down holds
\emph{anticipatory context} (waiting for results that may or may not
arrive). Escalation holds \emph{reactive context} (responding to
failures that have already occurred). Anticipatory context scales with
the number of delegations; reactive context scales with the failure rate.
As the system matures and lower levels improve, the failure rate
decreases, so escalation overhead \emph{decreases with capability}. In
hierarchical top-down, the number of delegations \emph{increases with
capability}: more capability means more subprocesses, each requiring a
delegation slot.

\textbf{Remark on convergence.} Hierarchical top-down that is
sufficiently deep, with intermediaries that handle most problems without
escalating, begins to \emph{resemble} escalation. The intermediary that
resolves a problem without reporting upward is functioning as an
escalation level: it holds no delegation context at the level above.
The convergence is one-directional: as hierarchical top-down becomes
more autonomous at lower levels, it converges toward escalation.
Escalation never converges toward top-down, because it never introduces
anticipatory delegation context.

\subsection{Event-driven ``top-down'' and causal direction}

A fourth objection: an event-driven top-down supervisor registers
callbacks, holds zero delegation state, and is invoked by subsystems on
completion or failure. This achieves escalation's $O(1)$ steady-state
cost while retaining ``top-down authority.''

The issue is not one of labeling --- it is not that we absorb every
efficient architecture under the name ``escalation'' to make the thesis
unfalsifiable. The issue is the \emph{causal direction of the signal}.
The structural distinction between top-down and escalation is:

\begin{itemize}
\item \textbf{Top-down}: the upper level initiates. The signal goes
   \emph{downward}. The upper level says ``do this'' and waits for the
   result.
\item \textbf{Escalation}: the lower level signals. The signal goes
   \emph{upward}. The lower level says ``I cannot handle this'' and the
   upper level engages reactively.
\end{itemize}

In the event-driven architecture, who initiates the callback? The
\emph{subsystem}. The subsystem detected the condition --- whether
completion or failure --- and sent the signal \emph{upward}. The upper
level did not poll, did not check, did not anticipate. It was invoked
from below. This is bottom-up signaling regardless of what the
architecture is called.

The claimed top-down authority --- the ability to initiate, cancel, or
redirect --- is unavailable to a supervisor holding zero state. A
supervisor with no delegation context does not know what is happening
below and cannot redirect what it is not tracking. When the callback
arrives, the supervisor must \emph{reload context} to interpret the
result --- exactly what an escalation level does when a failure signal
arrives.

The point is not definitional but empirical: if you trace the causal
arrow in the event-driven architecture, it points upward. The
architecture may have been \emph{designed} as top-down (the supervisor
originally assigned the task), but in steady-state operation, the
information flow is bottom-up. The label ``top-down'' describes the
initial assignment; the ongoing control structure is escalation.

\subsection{Hidden costs of escalation during learning}

A second objection: escalation achieves $O(1)$ steady-state cost, but
this assumes lower levels are competent. During learning --- when lower
levels frequently fail --- escalation produces frequent failure signals
at $L_n$, burdening it with escalation context.

This is correct but not fatal to the argument. During learning, both
architectures burden $L_n$: top-down burdened by delegation overhead
plus the novel problems themselves; escalation burdened only by the
novel problems (the failures that escalate). The difference: under
top-down, the overhead is \emph{additive} to the learning load
(delegation context plus problem context). Under escalation, the
overhead \emph{is} the learning load (failure context only). Escalation
during learning is not $O(1)$ --- it is $O(\text{failure rate} \cdot
|K^{[a]}|)$ --- but it never exceeds the intrinsic difficulty of the
problems themselves. Top-down adds coordination overhead on top of
intrinsic difficulty.

As learning proceeds and lower levels encode proven actions, escalation
load decreases monotonically. Top-down delegation overhead does not
decrease with learning --- it grows with the number of managed processes,
which increases as the system becomes more capable.

\subsection{Empirical dominance of dense attention}

A third objection: transformers with full self-attention --- the purest
top-down architecture, where every token attends to every other --- have
dominated alternative architectures for a decade. If top-down truly
degraded, why does dense attention win?

The answer is in the trajectory, not the snapshot. Dense attention
dominates \emph{today} at \emph{current scales}. But the trend is
toward escalation-like structure:

\begin{itemize}
\item \textbf{Sparse attention} (learned sparsity patterns, local
   windows, routing) reduces $O(n^2)$ by learning which edges to keep
   --- learned graph pruning \cite{McCarthy2026a}.
\item \textbf{Mixture-of-experts} routes tokens to specialized
   sub-networks rather than evaluating all parameters for all tokens ---
   escalation at the parameter level.
\item \textbf{Retrieval-augmented generation} retrieves relevant context
   from an external store rather than holding everything in the context
   window --- graph traversal from a query node.
\item \textbf{``Lost in the middle''} \cite{Liu2024}: empirical evidence
   that inference quality degrades as context fills with irrelevant
   content --- precisely the $\rho$ degradation of
   Theorem~\ref{thm:topdown}.
\end{itemize}

The industry is converging toward escalation. Dense attention is the
starting point, not the destination. The cost pressures the theorem
identifies are driving the convergence.

\textbf{Remark: the LLM theory of the case.}
Theorem~\ref{thm:topdown} does not claim that large language models
\emph{cannot work}. Dense attention over a large context window is a
viable strategy as long as the context can cover the problem space ---
that is, as long as the system operates below the crossover
(Section~\ref{sec:crossover}). Current LLMs succeed precisely because
their context windows are large enough, relative to the problems they
are asked to solve, that $\rho$ degradation is tolerable.

What the theorem establishes is that this strategy is
\emph{suboptimal and expensive}. Full self-attention recomputes the
pairwise adjacency structure from scratch at every query: $O(n^2)$ per
inference step. A system with stored adjacency pays $O(d)$ per
retrieval, where $d$ is the local degree. The cost ratio grows
quadratically with context size. For a 128k-token context window, the
system computes $\sim$$10^{10}$ pairwise attention weights per layer per
query, the vast majority of which are near-zero (irrelevant pairs). This
is the computational cost of not having a graph.

The theorem also predicts a ceiling. As the problems LLMs are asked to
solve grow in complexity --- requiring more knowledge, more reasoning
steps, more integration across domains --- the context required grows
with the problem, and $\rho$ degrades: ``lost in the middle''
\cite{Liu2024} is the empirical manifestation. Growing the context
window defers the crossover but does not eliminate it, because the
overhead of attending to irrelevant content grows quadratically while
the benefit of additional context grows sub-linearly (most added context
is irrelevant to any specific problem).

The LLM architecture is not invalidated by this result --- it is
\emph{located} by it. LLMs operate in the pre-crossover regime where
dense top-down attention is functional but costly. The industry
trajectory toward sparse attention, mixture-of-experts, and
retrieval-augmented generation is movement toward the post-crossover
architecture the theorem identifies as necessary: stored adjacency,
selective retrieval, escalation on failure.

\textbf{Remark on scope.} Theorem~\ref{thm:topdown} is not a claim
that top-down control is wrong in all cases. Occasional top-down
intervention --- where $L_n$ deliberately overrides a lower-level
response --- is a valid and sometimes necessary operation. The theorem
establishes that top-down allocation as the \emph{primary} control
architecture degrades under bounded $C_n$. The distinction is between
top-down as occasional intervention (compatible with escalation) and
top-down as the default routing mechanism (incompatible with bounded
context).

\textbf{Biological measurement.} The empirical signature of escalation
is the transition from active to automatic processing across skill
acquisition. When an informed action $a$ is first learned, it consumes
$L_n$ capacity: the learner must actively reason through each step. As
encoding proceeds --- as $K^{[a]}$ migrates from $L_n$ to lower levels
--- $a$ ceases to consume $L_n$ capacity. This is precisely the
escalation architecture: $a$ has been encoded at a level where it no
longer requires active context. The ubiquity of this transition across
motor learning, language acquisition, and expert skill development
\cite{Kahneman2011} is evidence that biological systems implement
escalation rather than persistent top-down allocation.

% ------------------------------------------------------------------
\section{Comparison: Escalation vs.\ Top-Down}
\label{sec:comparison}
% ------------------------------------------------------------------

The two architectures can be compared directly on their steady-state
context demands.

\subsection{Escalation architecture}

Under escalation, $L_n$ engages only when lower levels fail to bound
uncertainty for a given problem.

\begin{itemize}
\item \textbf{Steady-state context per problem:} $O(|K^{[a]}|)$.
   Only the knowledge relevant to the currently escalated problem
   occupies $L_n$. No delegation overhead is incurred because lower
   levels are not ``managed'' --- they operate autonomously.

\item \textbf{Effect of knowledge growth:} As $K$ grows and more
   actions are encoded at lower levels, \emph{fewer} problems reach
   $L_n$. The load on active context decreases or remains bounded as
   capability increases.

\item \textbf{Scaling property:} Capability growth \emph{relieves}
   $L_n$ rather than burdening it. Each successfully encoded action
   is one fewer demand on active context.
\end{itemize}

\subsection{Top-down allocation architecture}

Under top-down allocation, $L_n$ holds delegation context for all
managed subprocesses.

\begin{itemize}
\item \textbf{Steady-state context per problem:}
   $O(|K^{[a]}| + k \cdot \alpha)$, where $k$ is the number of
   concurrent delegations. The delegation overhead $k \cdot \alpha$
   is additive to the problem's own knowledge requirements.

\item \textbf{Effect of capability growth:} As $A$ grows, $k$ grows
   proportionally (the supervisor must track more subprocesses).
   $L_n$ load increases with capability.

\item \textbf{Scaling property:} Capability growth \emph{burdens}
   $L_n$. Each new capability adds to the supervisory overhead,
   eventually saturating active context.
\end{itemize}

\subsection{The inversion}

The escalation architecture inverts the scaling relationship between
capability and cognitive load:

\begin{center}
\begin{tabular}{@{}lcc@{}}
\toprule
& \textbf{Escalation} & \textbf{Top-Down} \\
\midrule
$L_n$ engagement & On failure only & Continuous supervision \\
Steady-state context & $O(|K^{[a]}|)$ & $O(|K^{[a]}| + k \cdot
\alpha)$ \\
As $|A|$ grows & $L_n$ load $\downarrow$ or bounded &
$L_n$ load $\uparrow$ proportionally \\
At capacity & Novel problems get full $C_n$ &
Novel problems get $C_n - k\alpha$ \\
\bottomrule
\end{tabular}
\end{center}

Under escalation, a system that has encoded 10{,}000 actions at lower
levels has \emph{more} active context available for novel problems than
when it had encoded 100 --- because 9{,}900 additional problems are
now handled without $L_n$ engagement. Under top-down allocation, a
system managing 10{,}000 subprocesses has \emph{less} active context
available than when it managed 100 --- because 9{,}900 additional
delegation slots consume context.

This is not a quantitative difference but a qualitative inversion of
the scaling direction. It explains why expertise feels like
\emph{effortlessness}: the expert has encoded actions that the novice
must actively reason through, freeing $L_n$ for genuinely novel
problems.

% ------------------------------------------------------------------
\section{Engagement with Prior Work}
\label{sec:prior}
% ------------------------------------------------------------------

\subsection{Free Energy Principle and Active Inference}

The Free Energy Principle \cite{Friston2010,Friston2017} posits that
biological systems minimize variational free energy, with attention
implemented as precision weighting: higher cortical levels set the gain
on prediction errors at lower levels. This is a top-down architecture
in the sense of Theorem~\ref{thm:topdown}: the supervisory level must
maintain a representation of expected precision for each lower-level
channel, and must evaluate incoming prediction errors against these
expectations.

Theorem~\ref{thm:topdown} shows that this precision allocation
degrades under bounded $C_n$ because the delegation overhead ---
maintaining expected precision for $k$ channels --- fills the active
context window proportionally to the number of monitored channels.
This is a critique of the architectural \emph{prescription}, not of the
minimization objective. Free energy minimization may well be the
correct objective function; the question is whether the control
architecture that implements it should be top-down (precision
allocation from above) or bottom-up (escalation on prediction error
that exceeds local capacity). Under bounded context, escalation
dominates.

More precisely: FEP's precision weighting may be optimal when $C_n$ is
not binding --- when active context is large relative to the number of
monitored processes. Under bounded context, the overhead of maintaining
precision weights for all channels eventually exceeds the benefit of
optimal allocation, and escalation (where channels signal only on
failure) achieves better inference quality at $L_n$.

\subsection{System M: Dupoux, LeCun, and Malik}

Dupoux, LeCun, and Malik \cite{DupouxLeCunMalik2026} propose a
three-component architecture: System~A (observation/world models),
System~B (action/planning), and System~M (meta-control), where M is a
fixed evolutionary routing policy that directs when A and B engage.
System~M monitors epistemic signals and routes data between
subsystems.

This is explicitly top-down: M sits above A and B, monitoring their
states and directing their activity. Theorem~\ref{thm:topdown}
establishes that M's monitoring itself consumes $C_n$ proportionally
to the number of subsystems and channels it oversees. As the
architecture scales --- more world models, more action policies, more
channels --- M's overhead grows, eventually crowding out the capacity
for the reasoning that M is supposed to coordinate.

Their architecture would benefit from replacing top-down routing with
escalation: subsystems A and B signal M only on failure (when local
uncertainty exceeds a threshold), rather than M continuously polling
all subsystems for state. Under escalation, M engages only when a
subsystem cannot resolve its own uncertainty --- and the overhead at
M's level is proportional to the \emph{failure rate}, not the
\emph{system size}.

\subsection{Global Workspace Theory}

Global Workspace Theory \cite{Baars1988,Dehaene2011} already
describes what is essentially an escalation architecture: specialized
local processors handle routine computation autonomously, and the
global workspace is engaged only when local processing fails or when
integration across modules is required. Broadcasting to the global
workspace is the cognitive analogue of escalation to $L_n$.

Theorem~\ref{thm:topdown} provides the formal grounding that GWT
assumes architecturally. GWT \emph{asserts} that local processors
handle routine work and the global workspace handles exceptions;
Theorem~\ref{thm:escalation} \emph{derives} the escalation criterion
from uncertainty-reduction cost, and Theorem~\ref{thm:topdown}
\emph{proves} why the alternative (top-down allocation from the global
workspace) fails. The $O(k \cdot \alpha)$ overhead is the formal
reason why global workspace engagement must be the exception rather
than the rule.

\subsection{System 1 and System 2}

Kahneman's System~1/System~2 framework \cite{Kahneman2011} is the
popular binary version of the encoding hierarchy. System~1 (fast,
automatic, low-cost) corresponds to lower levels of $E(c,r)$;
System~2 (slow, deliberate, high-cost) corresponds to $L_n$.

The continuous parameterization $E(c,r)$ corrects the binary
discretization. There are not two systems but a continuous space of
encodings with substrate-determined interfaces. A motor reflex
($c \approx 0$, $r \approx 0$) and a well-practiced heuristic
($c$ low, $r$ moderate) are both ``System~1'' in Kahneman's
taxonomy, but they occupy different regions of $E(c,r)$ with different
costs and different revisability. The binary collapses this structure
into a distinction that, while useful for popular exposition, obscures
the continuous nature of encoding and the multiplicity of interfaces.

Theorem~\ref{thm:escalation} also corrects the implicit control
assumption: System~2 does not ``decide'' when to override System~1.
Rather, System~1 (lower levels) handles problems autonomously, and
System~2 ($L_n$) engages when System~1 fails. The direction of
control is bottom-up, not top-down.

\subsection{Resource Rationality}

Griffiths, Lieder, and Goodman \cite{GriffithsEtAl2015} establish
that the optimal algorithm for a cognitive task depends on the
computational resources available --- resource rationality. The
optimal unconstrained algorithm may be suboptimal under resource
limitations.

Theorem~\ref{thm:topdown} is precisely an instance of this principle.
Top-down allocation may be the optimal control architecture when $C_n$
is unbounded: with infinite active context, the overhead of maintaining
delegation state for all subprocesses is negligible, and the benefit
of optimal routing is maximal. Under bounded $C_n$, the overhead is
no longer negligible, and escalation --- a ``simpler'' architecture
that does not attempt global optimization of routing --- dominates
because it conserves the scarce resource (active context) for genuine
inference rather than spending it on coordination.

Resource rationality provides the general principle;
Theorem~\ref{thm:topdown} provides the specific result for the
control architecture domain: the resource-rational control
architecture under bounded context is escalation.

% ------------------------------------------------------------------
\section{Conclusion}
\label{sec:conclusion}
% ------------------------------------------------------------------

We have shown that escalation dominance is a \emph{scaling result}.
Top-down allocation is viable --- and may be optimal --- when the action
repertoire $|A|$ fits within bounded context $C_n$. But intelligence
requires $|A| \gg C_n$: an entity must know far more than it can hold
in mind at once. Once $|A|$ exceeds the crossover point
$|A|^* = C_n / 2\alpha$ (Section~\ref{sec:crossover}), top-down
delegation overhead dominates active context and inference quality
degrades monotonically. The crossover is permanent under bounded $C_n$.

The argument has three components:

\begin{enumerate}
\item Theorem~\ref{thm:escalation} establishes that informed actions
   are handled at the lowest sufficient level, with escalation
   occurring only on failure.

\item Theorem~\ref{thm:topdown} establishes that top-down allocation
   degrades inference quality as $|A|$ grows, whether the architecture
   is flat ($k = |A|$) or hierarchical ($k = b \cdot \log_b |A|$).
   Recovery requires unbounded growth of $C_n$, which contradicts the
   boundedness assumption.

\item The crossover analysis (Section~\ref{sec:crossover}) identifies
   the regime boundary: top-down works below $|A|^*$, fails above it.
   Every system we recognize as intelligent operates well above it.
\end{enumerate}

The combination yields a precise result: escalation is not universally
superior to top-down --- it is superior at every scale that matters.
$L_n$ is the level of last resort, not the seat of control.
Capability growth, under escalation, \emph{relieves} the highest level
rather than burdening it. Under top-down allocation, capability growth
\emph{burdens} the highest level until it saturates.

This result has implications for both cognitive science and AI
architecture. In cognitive science, it provides the formal grounding
for Global Workspace Theory's architectural assumptions and corrects
the implicit top-down control assumption in System~1/System~2 accounts.
In AI architecture, it identifies a structural ceiling on centralized
routing: systems scaling beyond $|A|^*$ should implement escalation ---
subsystems that operate autonomously and signal on failure --- rather
than supervisory routing that attempts to manage all subprocesses from
above. The trajectory of transformer architectures, from dense
attention toward sparse attention, mixture-of-experts, and
retrieval-augmented generation, is convergence toward this result.

The series to this point has established two results: knowledge under
bounded context necessarily forms a graph \cite{McCarthy2026a}; and the
control architecture over that graph is necessarily escalation under
bounded context (this paper). Together, these results constrain the
architecture of any system that accumulates knowledge, acts under
uncertainty, and operates within bounded active context --- which is
every physical system that persists.

\bibliographystyle{plain}
\begin{thebibliography}{99}

\bibitem{Baars1988}
B.~J. Baars.
\newblock \emph{A Cognitive Theory of Consciousness}.
\newblock Cambridge University Press, 1988.

\bibitem{Clark2015}
A.~Clark.
\newblock \emph{Surfing Uncertainty: Prediction, Action, and the Embodied
  Mind}.
\newblock Oxford University Press, 2015.

\bibitem{Cowan2001}
N.~Cowan.
\newblock The magical number 4 in short-term memory: A reconsideration of
  mental storage capacity.
\newblock \emph{Behavioral and Brain Sciences}, 24(1):87--114, 2001.

\bibitem{Cover2006}
T.~M. Cover and J.~A. Thomas.
\newblock \emph{Elements of Information Theory}.
\newblock Wiley-Interscience, 2nd edition, 2006.

\bibitem{Dehaene2011}
S.~Dehaene.
\newblock \emph{Consciousness and the Brain: Deciphering How the Brain Codes
  Our Thoughts}.
\newblock Viking Press, 2011.

\bibitem{DupouxLeCunMalik2026}
E.~Dupoux, Y.~LeCun, and J.~Malik.
\newblock Why {AI} systems don't learn and what to do about it: Lessons on
  autonomous learning from cognitive science.
\newblock \emph{arXiv preprint arXiv:2603.15381}, 2026.

\bibitem{Friston2010}
K.~Friston.
\newblock The free-energy principle: A unified brain theory?
\newblock \emph{Nature Reviews Neuroscience}, 11(2):127--138, 2010.

\bibitem{Friston2017}
K.~Friston.
\newblock Active inference and learning.
\newblock \emph{Neuroscience \& Biobehavioral Reviews}, 68:862--879, 2017.

\bibitem{GriffithsEtAl2015}
T.~L. Griffiths, F.~Lieder, and N.~D. Goodman.
\newblock Rational use of cognitive resources: Levels of analysis between the
  computational and the algorithmic.
\newblock \emph{Topics in Cognitive Science}, 7(2):217--229, 2015.

\bibitem{Kahneman2011}
D.~Kahneman.
\newblock \emph{Thinking, Fast and Slow}.
\newblock Farrar, Straus and Giroux, 2011.

\bibitem{Liu2024}
N.~F. Liu, K.~Lin, J.~Hewitt, A.~Paranjape, M.~Bevilacqua, F.~Petroni, and
  P.~Liang.
\newblock Lost in the middle: How language models use long contexts.
\newblock \emph{Transactions of the Association for Computational Linguistics},
  12:157--173, 2024.

\bibitem{McCarthy2026a}
P.~D. McCarthy.
\newblock Graph structure is necessary for information preservation under
  bounded context.
\newblock Working paper, 2026.

\bibitem{Simon1955}
H.~A. Simon.
\newblock A behavioral model of rational choice.
\newblock \emph{The Quarterly Journal of Economics}, 69(1):99--118, 1955.

\bibitem{Simon1962}
H.~A. Simon.
\newblock The architecture of complexity.
\newblock \emph{Proceedings of the American Philosophical Society},
  106(6):467--482, 1962.

\bibitem{Williamson1967}
O.~E. Williamson.
\newblock Hierarchical control and optimum firm size.
\newblock \emph{Journal of Political Economy}, 75(2):123--138, 1967.

\end{thebibliography}

\end{document}
